\documentclass[11pt]{article}
\usepackage[margin=1in]{geometry}
\usepackage{amsmath,amssymb,amsthm,booktabs,array}
\usepackage[hidelinks]{hyperref}
\newtheorem{theorem}{Theorem}
\newtheorem{lemma}{Lemma}
\newtheorem{proposition}{Proposition}
\newcommand{\R}{\mathbb R}
\newcommand{\SOS}{\operatorname{SOS}}
\newcommand{\tr}{\operatorname{tr}}
\newcommand{\C}{\mathcal C}
\title{Nonconvexity of fixed odd-power sum-of-squares cones}
\author{Alec Kriebel\\\small Independent researcher\\
\small\href{https://orcid.org/0009-0001-9320-500X}{ORCID: 0009-0001-9320-500X}}
\date{6 October 2026}
\begin{document}
\maketitle
\begin{abstract}
For a fixed odd integer $q$, consider the nonnegative real homogeneous
forms whose $q$th powers are sums of squares. We give a negative answer to
the general-dimensional convexity question recorded by Reznick in 2023.
For sextics and $q=3$, the cone is nonconvex in the explicit finite dimension
$3\cdot10^{62}$. For every odd $q\ge3$, it is nonconvex in some finite
dimension depending on $q$. The proof combines a classical modified
Motzkin form with an SOS cube, separation from the SOS cone, positivity of
product truncated functionals, and elementary distinct-index counting.
An explicit rational moment functional gives the displayed dimension;
its $220\times220$ local positivity certificate can be checked without
floating-point arithmetic. No ternary or smallest-dimension conclusion is made.
\end{abstract}

\section{The question and the result}
Let $\R[x_1,\ldots,x_n]_m$ be the real forms of degree $m$, where $m$ is
positive and even. For a positive odd integer $q$, set
\[
 \C(n,m,q)=\{f\in\R[x_1,\ldots,x_n]_m:
                  f\ge0\text{ on }\R^n,\quad f^q\in\SOS\}.
\]
Here SOS means a finite sum of squares of real polynomials. Nonnegativity
is automatic from the second condition when $q$ is odd. These sets are
closed and stable under multiplication by nonnegative scalars; whether
they are convex is a separate question. Reznick asks this for general $n,m,q$
in his 2023 Oberwolfach contribution~\cite[p.779]{OWR}. The same question
appears in~\cite[\S6]{BKR}.

\begin{theorem}\label{thm:main}
Put $N=10^{62}$ and $n=3N$. On disjoint variable triples $(x_i,y_i,z_i)$,
define
\[
 p_i=x_i^4y_i^2+x_i^2y_i^4+z_i^6-x_i^2y_i^2z_i^2\qquad(1\le i\le N).
\]
Each $p_i$ belongs to $\C(n,6,3)$, whereas
$N^{-1}\sum_{i=1}^N p_i$ does not. Thus $\C(3\cdot10^{62},6,3)$ is not convex.
Moreover, for every fixed odd $q\ge3$ there is a finite $n(q)$ for which
$\C(n(q),6,q)$ is not convex.
\end{theorem}

The explicit dimension is a sufficient bound, with no optimality claim.
The theorem refutes the general convexity assertion; it does not settle
the ternary fixed-power question, classify parameter triples, or give a
common dimension for all odd exponents. Each counterexample persists on
adding unused variables, because restriction of an SOS to those variables
equal to zero is again SOS. The case $q=1$ and classical cases in which
every nonnegative form is SOS remain consistent with the theorem.

The ingredients are established. The exact modified Motzkin seed is due
to Berg--Christensen--Jensen~\cite[Lemma~1]{BCJ}; its SOS cube is recorded
in~\cite{OWR,BKR}. Closure and positive-functional separation are classical
as well~\cite[Theorems~3--4]{BCJ}. Product pseudoexpectation positivity is
explicit in Barak--Kelner--Steurer~\cite[Lemma A.5]{BKS}. Distinct-index
versus collision asymptotics appear, for example, in~\cite[Lemma~2.1]{AB}
and~\cite[Proposition~6.4]{BR}. We apply these tools to obstruction at one
\emph{prescribed} odd power and supply an explicit finite certificate.
The convexity of the union over odd powers in~\cite[Theorem~5.3]{BKR}
and~\cite[Theorem~2.2]{Stubborn} allows the exponent to increase and does
not imply fixed-power convexity.

\section{Disjoint-copy obstruction at a fixed odd power}
We write $\R[u]_{\le r}$ for polynomials of total degree at most $r$ in
a fixed finite variable block $u$.

\begin{lemma}\label{lem:product}
Let $p\in\R[u]_m$, let $q$ be positive and odd, and put $d=qm/2$.
Suppose $L:\R[u]_{\le2d}\to\R$ satisfies
\[
 L(1)=1,\qquad L(h^2)\ge0\ (\deg h\le d),\qquad a_1:=L(p)<0.
\]
For all sufficiently large finite integers $s$, the polynomial
$(p(u_1)+\cdots+p(u_s))^q$ on disjoint copies of $u$ is not SOS.
\end{lemma}
\begin{proof}
The local moment matrix
$M=(L(u^{\alpha+\beta}))_{|\alpha|,|\beta|\le d}$ is positive semidefinite.
On polynomials of degree at most $2d$ in \emph{each} of the $s$ blocks,
define a product functional on monomials by
\[
 L_s\!\left(\prod_{i=1}^s u_i^{\alpha_i}\right)
                  =\prod_{i=1}^s L(u^{\alpha_i}).
\]
The matrix controlling $L_s(H^2)$ for all $H$ of \emph{total} degree at
most $d$ is a principal submatrix of $M^{\otimes s}$, hence is positive
semidefinite. Consequently $L_s$ is nonnegative on every sum of such
squares. This is the standard product-positivity construction~\cite{BKS}.
All spaces and products here are finite; no representing measure is assumed.

Put $a_j=L(p^j)$ for $0\le j\le q$, so $a_0=1$, and let
$S_s=\sum_{i=1}^s p(u_i)$. Expansion by partitions of the $q$ labelled
factors gives
\begin{equation}\label{eq:partitions}
 L_s(S_s^q)=\sum_{\pi\in\operatorname{Part}([q])}
              (s)_{|\pi|}\prod_{B\in\pi}a_{|B|},
 \qquad (s)_r=s(s-1)\cdots(s-r+1).
\end{equation}
For each partition, assigning its distinct blocks to distinct copies
has $(s)_{|\pi|}$ choices. The unique partition with $q$ blocks is the
singleton partition. Thus the expression is a polynomial in $s$ of degree
$q$, with leading coefficient $a_1^q<0$; all collisions contribute lower
degree. It is strictly negative for all sufficiently large integers $s$.

An SOS polynomial of degree at most $2d$ can involve only square summands
of degree at most $d$: otherwise its highest-degree homogeneous part
would be a nonzero sum of squares, and could not cancel. Therefore the
negative value excludes every SOS representation of $S_s^q$.
\end{proof}

\begin{lemma}\label{lem:separation}
If a form $p$ is not SOS and $q$ is positive and odd, then with $d=qm/2$
there exists an $L$ satisfying the hypotheses of Lemma~\ref{lem:product}.
\end{lemma}
\begin{proof}
The cone of SOS polynomials of degree at most $2d$ is closed. Indeed, let
$v$ be the vector of all monomials of degree at most $d$ and write a
coefficient-convergent SOS sequence as $v^TQ_jv$, $Q_j\succeq0$.
For a standard Gaussian integral $\Gamma$,
$G=\Gamma(vv^T)$ is positive definite. Coefficient convergence bounds
$\Gamma(v^TQ_jv)=\tr(Q_jG)$, which bounds $\tr Q_j$, and hence the matrices
$Q_j$. A convergent subsequence yields a positive semidefinite Gram
matrix for the limit.

The same highest-degree argument used above shows that a non-SOS form
$p$ is outside this larger finite-degree SOS cone. Finite-dimensional
separation from a closed convex cone gives a functional $L_0$ nonnegative
on that cone with $L_0(p)<0$. For sufficiently small $\varepsilon>0$,
$L_0+\varepsilon\Gamma$ still has negative value on $p$ and has positive
definite moment matrix. Its value on $1$ is positive. Dividing by that
value gives the required $L$. This is classical SOS duality~\cite{BCJ}.
\end{proof}

The power map $f\mapsto f^q$ is continuous in the coefficient topology.
The closedness just proved therefore also gives closedness of each fixed-power
set $\C(n,m,q)$ as the inverse image of the fixed-degree SOS cone.

In particular, if a non-SOS nonnegative form $p$ has an SOS $q$th power,
its disjoint copies belong to the same ambient $\C$, whereas a sufficiently
large finite average does not. Convexity implies closure under every
finite convex combination by induction, so that average refutes convexity.
If a violating pair is desired, take the least failing prefix of the
finite sum in the same ambient variable space; its preceding prefix and
next summand are inside the cone. This asserts existence of a pair,
without assigning an SOS power to any untested prescribed prefix.

\section{The classical seed}
Set
\begin{equation}\label{eq:seed}
 p=x^4y^2+x^2y^4+z^6-x^2y^2z^2.
\end{equation}
AM--GM gives $x^4y^2+x^2y^4+z^6\ge3x^2y^2z^2$, hence $p\ge0$.
The dehomogenization $p(x,y,1)$ is exactly the example in~\cite[Lemma~1]{BCJ}.
Below we also independently certify $p\notin\SOS$ by a positive-on-squares
functional with $L(p)=-1$.
The credited SOS identity for $p^3$~\cite{OWR,BKR} is displayed in
Appendix~\ref{app:cube} and checked by exact polynomial expansion in the
supplement. For every odd $q\ge3$,
\[
                 p^q=p^3\bigl(p^{(q-3)/2}\bigr)^2\in\SOS.
\]
Lemmas~\ref{lem:product} and~\ref{lem:separation} therefore prove the second
assertion of Theorem~\ref{thm:main}. The functional may depend on $q$;
no degree 18 functional is reused outside its domain.

\section{An explicit rational functional for the cube}
For $\alpha\in\mathbb N^3$, let $g_\alpha$ be the standard Gaussian moment:
it is zero if any coordinate is odd, and otherwise
$g_\alpha=\prod_{j=1}^3(\alpha_j-1)!!$, with $(-1)!!=1$.
We define $\mu_\alpha=L(x^{\alpha_1}y^{\alpha_2}z^{\alpha_3})$
for $|\alpha|\le18$. Set all moments with an odd coordinate to zero.
For even coordinates, put $\mu_0=1$, and
$\mu_\alpha=g_\alpha/1024$ when $|\alpha|=2$ or $4$.
For degree 6, write $\alpha=2(a,b,c)$ and use
\[
\begin{array}{c|rrrrr}
(a,b,c)&(3,0,0)&(0,3,0)&(2,1,0)&(1,2,0)&(0,0,3)\\\hline
\mu_{2(a,b,c)}&2^{24}&2^{24}&1&1&1\\[2pt]
(a,b,c)&(1,1,1)&(1,0,2)&(0,1,2)&(2,0,1)&(0,2,1)\\\hline
\mu_{2(a,b,c)}&4&32&32&4096&4096
\end{array}
\]
For $|\alpha|=2r$, $4\le r\le9$, put $\mu_\alpha=T_rg_\alpha$, where
\begin{equation}\label{eq:scales}
\begin{array}{c|rrrrrr}
r&4&5&6&7&8&9\\\hline
T_r&10^{16}&10^{34}&10^{53}&10^{73}&10^{94}&10^{116}.
\end{array}
\end{equation}
This table specifies the entire order-nine moment matrix, with
$\binom{12}{3}=220$ monomials. Coordinate parity splits it into eight
blocks of size at most 35.

\begin{proposition}\label{prop:certificate}
The functional specified above is positive on every nonzero square of
degree at most 18: $L(h^2)>0$ for $0\ne h$, $\deg h\le9$.
\end{proposition}
\begin{proof}
Here is a compressed exact rational certificate and its verification rule.
Order multi-indices lexicographically by increasing first coordinate,
then second, then third, within the set $|\alpha|\le r$; retain this order
in each parity block. At order 3 the leading principal minors are as follows.
The entries in each row list all such minors for that block.
\begin{center}\small
\begin{tabular}{c l}
\toprule
Parity&Leading principal minors at order 3\\\midrule
000&$1,\ 3071/1048576,\ 2047/268435456,\ 5117/274877906944$\\
001&$1/1024,\ 1015/1048576,\ 3109055/1048576,\ 1056109395/131072$\\
010&$1/1024,\ 32767/1048576,\ 16642474743/32768,\ 257328824345/1048576$\\
100&$1/1024,\ 32767/1048576,\ 16359/1048576,\ 257328824345/1048576$\\
011,101,110&$1/1024$ each\\
111&$4$\\\bottomrule
\end{tabular}
\end{center}
They are positive, so Sylvester's criterion establishes order 3 positivity.
For $r=4,\ldots,9$, split each block with new degree-$r$ monomials into
old and new rows. It has the form
\[
 \begin{pmatrix}A&B\\B^T&T_rG\end{pmatrix},\qquad
 G=(g_{\alpha+\beta})_{\substack{|\alpha|=|\beta|=r\\
                            \alpha,\beta\text{ of the given parity}}}.
\]
The Gaussian matrix $G$ is positive definite. The lower-degree block $A$
is positive definite by induction; $B$ uses already specified moments.
Let $C=B^TA^{-1}B$. Exact rational calculations give
\begin{equation}\label{eq:traces}
                   0\le\tr(G^{-1}C)<T_r
\end{equation}
for each of the four extending blocks at each order. The supplement's rational witness file
\path{TENSOR_MOMENT_CERTIFICATE.json} lists all 24 exact fractions,
the old/new dimensions and strict bounds. Its standalone checker rebuilds
$A,B,G$ from the displayed moments, solves the rational systems and
checks these fractions and inequalities; supplied positivity flags are
not used. It also independently computes all 220 positive pivots of the
full local blocks and verifies the exact LDL reconstruction.

Since $C\succeq0$, the largest eigenvalue of
$G^{-1/2}CG^{-1/2}$ is at most its trace. Equation~\eqref{eq:traces}
therefore implies $T_rG-C\succ0$. The Schur complement makes every
enlarged block positive definite. Blocks with no new rows are unchanged.
Induction proves the proposition. The finite computations in this step
are exact rational certificate checks, with no numerical eigensolver or SDP.
\end{proof}

\section{The finite-dimensional counterexample}
Exact expansion using the moment table yields
\begin{equation}\label{eq:values}
 a_1=L(p)=-1,\quad a_2=L(p^2)=11292\cdot10^{53},\quad
 a_3=L(p^3)=35039520\cdot10^{116}.
\end{equation}
The first value is $1+1+1-4$. For the other two, homogeneity means all
moments have degree 12 or 18, respectively, so these are Gaussian polynomial
moments times $T_6$ or $T_9$. The checker expands both polynomials exactly.
Proposition~\ref{prop:certificate} and $L(p)=-1$ also prove that $p$ is not SOS.

For $S_N=\sum_{i=1}^N p_i$, the cubic case of~\eqref{eq:partitions} is
\[
 L_N(S_N^3)=Na_3+3N(N-1)a_2a_1+N(N-1)(N-2)a_1^3.
\]
Since $a_1=-1$ and $a_2\ge1$, we obtain
\[
 L_N(S_N^3)\le N\{a_3-(N^2-1)\}<0,
\]
where the strict inequality follows from
\[
 35039520\cdot10^{116}<10^{124}-1=N^2-1\qquad(N=10^{62}).
\]
Lemma~\ref{lem:product}'s product positivity excludes every SOS for
$S_N^3$, and positive scaling excludes one for $(S_N/N)^3$.
Each $p_i^3$ has the credited SOS identity. Thus the average is a finite
convex combination of elements of $\C(3N,6,3)$ outside that set, proving
the first assertion of Theorem~\ref{thm:main}.

The functional $L$ is not a measure on real points: its negative value on
the nonnegative $p$ prevents that interpretation. Only positivity on the
required squares is asserted. The argument neither constructs an enormous
matrix numerically nor assumes an infinite product probability measure.

\section{Verification, priority and disclosure}
The accompanying verification archive contains this standalone source,
the complete rational certificate, a Python-standard-library checker,
expected deterministic output, hash/provenance records, the claim scope
and a dated priority note. The checker uses integers and rational fractions
to verify the local moment matrices, all initial minors and Schur traces,
the credited cube identity, the three seed moments and the exact finite
cubic sign. These computations support the finite certificate; the
arbitrary finite product argument and the every-odd-$q$ assertion are
proved analytically above. No formal proof-assistant certificate is claimed.

Two independent extensive bounded priority audits identified no exact
earlier fixed-power result in the sources examined. Their source versions,
reading extents and remaining gaps are disclosed in the supporting note.
The full final Papp--Alizadeh 2013 article and Schm\"udgen 1979 body were
unavailable; a captured Choi--Lam--Reznick 1995 scan was unread and unused.
Classical tensor/separation ingredients have independently accessible
primary proofs. Bibliographies, citation indexes and talk coverage have
limits. This is not a worldwide-firstness guarantee or a certificate of
present openness. The claim concerns the general question recorded in 2023.

AI tools were used extensively in derivation, literature comparison,
reproduction, writing and adversarial verification. This is an unrefereed
preprint and has not received human peer review. Mathematical and source
claims remain open to independent correction. The author metadata does
not imply affiliation with a research institution.

\appendix
\section{The credited cube identity}\label{app:cube}
For the seed~\eqref{eq:seed}, put
\begin{align*}
h_1&=x^5y^4-x^3y^4z^2,& h_2&=x^4y^5-x^4y^3z^2,\\
h_3&=x^4y^2z^3-x^2y^2z^5,&h_4&=x^2y^4z^3-x^2y^2z^5,\\
h_5&=xy^2z^6-x^3y^4z^2,&h_6&=x^2yz^6-x^4y^3z^2,\\
h_7&=x^2y^4z^3-x^4y^2z^3,&h_8&=x^4y^5-x^2yz^6,\\
h_9&=x^5y^4-xy^2z^6.&&
\end{align*}
Then the identity in~\cite{BKR}, specialized to this seed, is
\begin{align*}
p^3={}&\frac32\sum_{j=1}^9h_j^2+(z^9-2x^2y^2z^5)^2
 +(xyz^7-2x^3y^3z^3)^2\\
&+(x^3y^6-2x^3y^4z^2)^2+(x^3y^5z-2x^3y^3z^3)^2\\
&+(x^6y^3-2x^4y^3z^2)^2+(x^5y^3z-2x^3y^3z^3)^2
 +2(x^3y^3z^3)^2.
\end{align*}
Every square factor has degree 9 and every weight is positive.
This is the modified coefficient-$1$ seed, not the coefficient-$3$
original Motzkin form.

\begin{thebibliography}{9}
\raggedright
\bibitem{OWR} B. Reznick, The Odd Powers of the Motzkin Polynomial, etc.,
in \emph{Real Algebraic Geometry with a View toward Koopman
Operator Methods}, Oberwolfach Reports 14/2023, pp.~778--781.
\href{https://doi.org/10.4171/OWR/2023/14}{doi:10.4171/OWR/2023/14}.
\bibitem{BKR} G. Blekherman, K. Kozhasov and B. Reznick,
On odd powers of nonnegative polynomials that are not sums of squares,
\emph{Forum of Mathematics, Sigma} 14 (2026), e65.
\href{https://doi.org/10.1017/fms.2026.10221}{doi:10.1017/fms.2026.10221}.
\bibitem{BCJ} C. Berg, J. P. R. Christensen and C. U. Jensen,
A remark on the multidimensional moment problem,
\emph{Mathematische Annalen} 243 (1979), 163--169.
\href{https://doi.org/10.1007/BF01420423}{doi:10.1007/BF01420423}.
\bibitem{BKS} B. Barak, J. Kelner and D. Steurer,
Rounding Sum-of-Squares Relaxations, ECCC Report TR13-184 (2013),
Appendix A, Lemma A.5.
\url{https://eccc.weizmann.ac.il/report/2013/184/}.
\bibitem{AB} J. Acevedo and G. Blekherman,
Power Mean Inequalities and Sums of Squares, arXiv:2303.11823v1 (2023).
\url{https://arxiv.org/abs/2303.11823}.
\bibitem{BR} G. Blekherman and C. Riener,
Symmetric Nonnegative Forms and Sums of Squares,
arXiv:1205.3102v4 (2020), Proposition~6.4.
\url{https://arxiv.org/abs/1205.3102}.
\bibitem{Stubborn} L. Baldi, G. Blekherman, K. Kozhasov, D. Plaumann,
B. Reznick and R. Sinn, Stubborn Polynomials,
arXiv:2602.01191v1 (2026).
\url{https://arxiv.org/abs/2602.01191}.
\end{thebibliography}
\end{document}
